\documentclass[../../../include/open-logic-section]{subfiles}

\begin{document}

\olfileid[hi]{sth}{z}{power}
\olsection{घात समुच्चय}

हम एक और अभिगृहीत के साथ आगे बढ़ेंगे:

\begin{axiom}[घात समुच्चय]
किसी भी समुच्चय $A$ के लिए समुच्चय
$\Pow{A} = \Setabs{x}{x \subseteq A}$ का अस्तित्व है।
\[
	\forall A \exists P \forall x(x \in P \liff (\forall z \in x)z \in A)
\]
\end{axiom}

इसका औचित्य काफ़ी सीधा है। मान लें कि $A$ चरण $S$ पर बना है। तब $A$ के
सभी सदस्य $S$ से पहले उपलब्ध थे (\stagesacc{} से)। अतः पृथक्करण के औचित्य
में दिए तर्क की तरह $A$ का प्रत्येक उपसमुच्चय चरण $S$ तक बन चुका है। इसलिए
$S$ के बाद किसी भी चरण पर वे सभी एक समुच्चय में बनाए जाने के लिए उपलब्ध हैं।
और हम जानते हैं कि ऐसा कोई चरण है, क्योंकि $S$ अंतिम चरण नहीं है
(\stagessucc{} से)। अतः $\Pow{A}$ का अस्तित्व है।

घात समुच्चय का एक अच्छा परिणाम यह है:

\begin{prop}\label{thm:Products}
किन्हीं समुच्चयों $A, B$ का कार्तीय गुणनफल $A \times B$ अस्तित्व में है।
\end{prop}

\begin{proof}
घात समुच्चय और \olref[sth][z][pairs]{prop:pairsconsequences} से समुच्चय
$\Pow{\Pow{A \cup B}}$ का अस्तित्व है। अतः पृथक्करण से इस समुच्चय का
अस्तित्व है:
\[
	C = \Setabs{z \in \Pow{\Pow{A \cup B}}}{(\exists x \in A)(\exists y \in B) z = \tuple{x, y}}.
\]
अब किन्हीं $x \in A$ और $y \in B$ के लिए
\olref[sth][z][pairs]{prop:pairsconsequences} से समुच्चय $\tuple{x, y}$
का अस्तित्व है। इसके अतिरिक्त, चूँकि $x, y \in A \cup B$, हमारे पास
$\{x\}, \{x, y\} \in \Pow{A \cup B}$ और
$\tuple{x,y} \in \Pow{\Pow{A \cup B}}$ हैं। अतः $A \times B = C$।
\end{proof}

इस प्रमाण में घात समुच्चय पृथक्करण के साथ काम करता है। यह आश्चर्यजनक नहीं।
पृथक्करण के बिना घात समुच्चय बहुत \emph{शक्तिशाली} सिद्धांत न होता। आखिर
पृथक्करण हमें बताता है कि किसी समुच्चय के कौन-से उपसमुच्चय अस्तित्व रखते हैं
और इस प्रकार निर्धारित करता है कि प्रत्येक घात समुच्चय कितना ``भरा हुआ'' है।

\begin{prob}
दिखाइए कि किन्हीं समुच्चयों $A, B$ के लिए: (i) प्रांत $A$ और परास $B$
वाले सभी संबंधों के समुच्चय का अस्तित्व है; और (ii) $A$ से $B$ में सभी
फलनों के समुच्चय का अस्तित्व है।
\end{prob}

\begin{prob}
$A$ को समुच्चय और $\sim$ को $A$ पर कोई तुल्यता संबंध मानें। सिद्ध कीजिए कि
$A$ पर $\sim$ के अधीन तुल्यता वर्गों का समुच्चय, अर्थात
$\equivclass{A}{\sim}$, अस्तित्व में है।
\end{prob}

\end{document}
